% GENERATED FILE. Edit canonical lessons, then run npm run generate:books.
% canonical-source-hash: fnv1a64:387861629ea72eae
% canonical-lessons: RU-C389-zapreshchat, RU-C389-zapret, RU-C389-razreshat, RU-C389-razreshenie, RU-C389-pozvolyat

\chapter{Allowed and forbidden}
\label{ch:ru-allowed-and-forbidden}
\hlchaptermodality{389}

\begin{chapteropening}
\textbf{By the end of this chapter:} \emph{I can say what is forbidden and what is allowed, talk about a ban and a permit, and say what something lets me do.}

Complete the last lesson of chapter 389: I can say what is forbidden and what is allowed, talk about a ban and a permit, and say what something lets me do.
\end{chapteropening}

\section[\texorpdfstring{\ru{запрещать} (zapreshchát)}{запрещать (zapreshchát)}]{\ru{запрещать} (zapreshchát) — to forbid, to ban}

\label{lesson:RU-C389-zapreshchat}



\begin{quote}
Before the new one: say the Russian for to break, then the Russian for a violation.
\end{quote}

\subsection*{You'll want to know: \ru{запрещать}}
\textbf{\ru{запрещать}} (\emph{zapreshchát}) — \textquotedblleft{}to forbid, to ban\textquotedblright{}.

\begin{quote}
\ru{Врач} \ru{запрещает} \ru{мне} \ru{есть} \ru{сладкое}.
\end{quote}

\begin{quote}
\emph{Vrach zapreshcháyet mne yest sládkoye.}
\end{quote}

\begin{quote}
The doctor forbids me to eat sweet things.
\end{quote}

Imperfective; the perfective is \emph{\ru{запретить}} (\emph{\ru{запрещу}, \ru{запретишь}}). \emph{\ru{Запрещать} \ru{кому}} + infinitive. \emph{\ru{Запрещено}} — it is forbidden.

\begin{cousinweb}
\textbf{\ru{запрещать}} is built from \emph{\ru{за}-} and the root of \emph{\ru{претить}} \textquotedblleft{}to be repugnant\textquotedblright{}, once \textquotedblleft{}to threaten\textquotedblright{}; the \emph{\ru{щ}} is the Church Slavonic softening of \emph{\ru{т}}, where Russian has \emph{\ru{ч}}.
\end{cousinweb}

\begin{grammarlens}[title={What you've built}]
One more word for telling, explaining and describing at B1.
\end{grammarlens}

\subsection*{Your turn}
\begin{itemize}
\raggedright
  \item \emph{Say it:} \emph{zapreshchát}
  \item \emph{Say it:} \emph{zapreshchát}, once more
  \item \emph{Say it:} a sentence of your own with \emph{zapreshchát}
  \item \emph{Recall:} say the Russian for to break, then the Russian for a violation, then say \emph{zapreshchát} again
\end{itemize}

\begin{tcolorbox}[breakable,colback=teal!4,colframe=teal!35!black,title={Before you move on}]
What does \ru{запрещать} mean? (\textquotedblleft{}To forbid, to ban\textquotedblright{}.) Say it once more.
\end{tcolorbox}
\section[\texorpdfstring{\ru{запрет} (zaprét)}{запрет (zaprét)}]{\ru{запрет} (zaprét) — a ban}

\label{lesson:RU-C389-zapret}



\begin{quote}
Before the new one: say the Russian for a violation, then the Russian for to forbid.
\end{quote}

\subsection*{You'll want to know: \ru{запрет}}
\textbf{\ru{запрет}} (\emph{zaprét}) — \textquotedblleft{}a ban\textquotedblright{}.

\begin{quote}
\ru{В} \ru{нашем} \ru{городе} \ru{есть} \ru{запрет} \ru{на} \ru{курение} \ru{в} \ru{парках}.
\end{quote}

\begin{quote}
\emph{V náshem górode yest zaprét na kuréniye v párkakh.}
\end{quote}

\begin{quote}
In our town there is a ban on smoking in parks.
\end{quote}

Masculine. \emph{\ru{Запрет} \ru{на}} + accusative — a ban on.

\begin{cousinweb}
\textbf{\ru{запрет}} is formed from \emph{\ru{запретить}} \textquotedblleft{}to forbid\textquotedblright{}, built from \emph{\ru{за}-} and the root of \emph{\ru{претить}}.
\end{cousinweb}

\begin{grammarlens}[title={What you've built}]
One more word for telling, explaining and describing at B1.
\end{grammarlens}

\subsection*{Your turn}
\begin{itemize}
\raggedright
  \item \emph{Say it:} \emph{zaprét}
  \item \emph{Say it:} \emph{zaprét}, once more
  \item \emph{Say it:} a sentence of your own with \emph{zaprét}
  \item \emph{Recall:} say the Russian for a violation, then the Russian for to forbid, then say \emph{zaprét} again
\end{itemize}

\begin{tcolorbox}[breakable,colback=teal!4,colframe=teal!35!black,title={Before you move on}]
What does \ru{запрет} mean? (\textquotedblleft{}A ban\textquotedblright{}.) Say it once more.
\end{tcolorbox}
\section[\texorpdfstring{\ru{разрешать} (razreshát)}{разрешать (razreshát)}]{\ru{разрешать} (razreshát) — to allow, to permit}

\label{lesson:RU-C389-razreshat}



\begin{quote}
Before the new one: say the Russian for to forbid, then the Russian for a ban.
\end{quote}

\subsection*{You'll want to know: \ru{разрешать}}
\textbf{\ru{разрешать}} (\emph{razreshát}) — \textquotedblleft{}to allow, to permit\textquotedblright{}.

\begin{quote}
\ru{Мама} \ru{разрешает} \ru{мне} \ru{гулять} \ru{до} \ru{девяти}.
\end{quote}

\begin{quote}
\emph{Máma razresháyet mne gulyát do devyatí.}
\end{quote}

\begin{quote}
Mum lets me stay out until nine.
\end{quote}

Imperfective; perfective \emph{\ru{разрешить}}. \emph{\ru{Разрешать} \ru{кому}} + infinitive. It also means to resolve: \emph{\ru{разрешить} \ru{спор}}.

\begin{cousinweb}
\textbf{\ru{разрешать}} is built from \emph{\ru{раз}-} and \emph{\ru{решить}}, from Proto-Slavic \emph{*rěšiti} \textquotedblleft{}to untie, to loosen\textquotedblright{}: to allow is to untie a restriction.
\end{cousinweb}

\begin{grammarlens}[title={What you've built}]
One more word for telling, explaining and describing at B1.
\end{grammarlens}

\subsection*{Your turn}
\begin{itemize}
\raggedright
  \item \emph{Say it:} \emph{razreshát}
  \item \emph{Say it:} \emph{razreshát}, once more
  \item \emph{Say it:} a sentence of your own with \emph{razreshát}
  \item \emph{Recall:} say the Russian for to forbid, then the Russian for a ban, then say \emph{razreshát} again
\end{itemize}

\begin{tcolorbox}[breakable,colback=teal!4,colframe=teal!35!black,title={Before you move on}]
What does \ru{разрешать} mean? (\textquotedblleft{}To allow, to permit\textquotedblright{}.) Say it once more.
\end{tcolorbox}
\section[\texorpdfstring{\ru{разрешение} (razreshéniye)}{разрешение (razreshéniye)}]{\ru{разрешение} (razreshéniye) — permission, a permit}

\label{lesson:RU-C389-razreshenie}



\begin{quote}
Before the new one: say the Russian for a ban, then the Russian for to allow.
\end{quote}

\subsection*{You'll want to know: \ru{разрешение}}
\textbf{\ru{разрешение}} (\emph{razreshéniye}) — \textquotedblleft{}permission, a permit\textquotedblright{}.

\begin{quote}
\ru{Для} \ru{этого} \ru{нужно} \ru{разрешение}.
\end{quote}

\begin{quote}
\emph{Dlya étogo núzhno razreshéniye.}
\end{quote}

\begin{quote}
You need permission for this.
\end{quote}

Neuter. \emph{\ru{Разрешение} \ru{на} \ru{работу}} — a work permit. It is also the resolution of a screen.

\begin{cousinweb}
\textbf{\ru{разрешение}} is built on \emph{\ru{разрешить}} \textquotedblleft{}to allow\textquotedblright{} with the noun suffix \emph{-\ru{ение}}; the root is Proto-Slavic \emph{*rěšiti} \textquotedblleft{}to untie\textquotedblright{}.
\end{cousinweb}

\begin{grammarlens}[title={What you've built}]
One more word for telling, explaining and describing at B1.
\end{grammarlens}

\subsection*{Your turn}
\begin{itemize}
\raggedright
  \item \emph{Say it:} \emph{razreshéniye}
  \item \emph{Say it:} \emph{razreshéniye}, once more
  \item \emph{Say it:} a sentence of your own with \emph{razreshéniye}
  \item \emph{Recall:} say the Russian for a ban, then the Russian for to allow, then say \emph{razreshéniye} again
\end{itemize}

\begin{tcolorbox}[breakable,colback=teal!4,colframe=teal!35!black,title={Before you move on}]
What does \ru{разрешение} mean? (\textquotedblleft{}Permission, a permit\textquotedblright{}.) Say it once more.
\end{tcolorbox}
\section[\texorpdfstring{\ru{позволять} (pozvolyát)}{позволять (pozvolyát)}]{\ru{позволять} (pozvolyát) — to let, to allow; to afford}

\label{lesson:RU-C389-pozvolyat}



\begin{quote}
Before the new one: say the Russian for to allow, then the Russian for permission.
\end{quote}

\subsection*{You'll want to know: \ru{позволять}}
\textbf{\ru{позволять}} (\emph{pozvolyát}) — \textquotedblleft{}to let, to allow; to afford\textquotedblright{}.

\begin{quote}
\ru{Зарплата} \ru{не} \ru{позволяет} \ru{нам} \ru{купить} \ru{машину}.
\end{quote}

\begin{quote}
\emph{Zarpláta ne pozvolyáyet nam kupít mashínu.}
\end{quote}

\begin{quote}
Our pay does not let us buy a car.
\end{quote}

Imperfective; perfective \emph{\ru{позволить}}. \emph{\ru{Позволять} \ru{себе}} — to afford: \emph{\ru{я} \ru{не} \ru{могу} \ru{себе} \ru{это} \ru{позволить}}.

\begin{cousinweb}
\textbf{\ru{позволять}} is built from \emph{\ru{по}-} and \emph{\ru{воля}} \textquotedblleft{}will, freedom\textquotedblright{}, from Proto-Slavic \emph{*voľa} (EtymDB): to allow someone is to give them their will.
\end{cousinweb}

\begin{grammarlens}[title={What you've built}]
One more word for telling, explaining and describing at B1.
\end{grammarlens}

\subsection*{Your turn}
\begin{itemize}
\raggedright
  \item \emph{Say it:} \emph{pozvolyát}
  \item \emph{Say it:} \emph{pozvolyát}, once more
  \item \emph{Say it:} a sentence of your own with \emph{pozvolyát}
  \item \emph{Recall:} say the Russian for to allow, then the Russian for permission, then say \emph{pozvolyát} again
\end{itemize}

\begin{tcolorbox}[breakable,colback=teal!4,colframe=teal!35!black,title={Before you move on}]
What does \ru{позволять} mean? (\textquotedblleft{}To let, to allow; to afford\textquotedblright{}.) Say it once more.
\end{tcolorbox}
